% !PRINTED 168
\weberpage{168}

Wenn nun $w$ irgend eine unzerlegbare Function ist, die
weder in allen $u_r$ noch in allen $v_s$ aufgeht, so wählen wir in (3)
$r$ und $s$ so, dass $u_r$ das erste nicht durch $w$ theilbare $u$ ist, also
$u_{r-1}, u_{r-2}\dots$ durch $w$ theilbar sind, und dass ebenso $v_s$ das
erste, nicht durch $w$ theilbare $v$ wird. Dann kann auch $U_{r+s}$
nicht durch $w$ theilbar sein, weil alle Glieder, mit Ausnahme
des ersten, durch $w$ theilbar sind, d. h. $F(t)$ ist primitiv.

Daraus folgt unmittelbar, wenn $p$ und $q$ irgend welche ganze
rationale Functionen der $x$ sind, dass $pq$ der Theiler des Productes
der beiden imprimitiven Functionen $pf(t)$, $q\varphi(t)$ ist, also
der zweite Theil des Satzes V. Daraus folgt weiter:

\gesperrt{VI. Wenn eine ganze rationale Function. $F(t)$ der $n+1$ Variablen $x$ und $t$ in zwei Factoren zerlegbar ist, die in Bezug auf $t$ ganz, in Bezug auf die $x$ wenigstens rational sind, so ist sie auch in zwei Functionen zerlegbar, die in $x$ und $t$ ganz und rational sind.} % [sic] printed period after "Function" before F(t)

Denn nach der Voraussetzung giebt es eine ganze rationale
Function $w$ der $x$ allein und zwei ganze rationale Functionen
der $x$ und $t$, $f_1(t)$, $\varphi_1(t)$, so dass
\[
w\,F(t) = f_1(t)\,\varphi_1(t),
\]
und nach (V) muss $w$ in dem Product der Theiler von $f_1(t)$ und
$\varphi_1(t)$ aufgehen, so dass wir auch
\[
F(t) = f(t)\,\varphi(t)
\]
erhalten, worin $f(t)$, $\varphi(t)$ gleichfalls ganze rationale Functionen
von $x$ und $t$, und zwar in Bezug auf $t$ von demselben Grade wie
$f_1(t)$ und $\varphi_1(t)$ sind, w. z. b. w. --- Hieraus schliessen wir weiter,
dass zwei Functionen $F(t)$, $f(t)$, die als ganze rationale Functionen
der $n+1$ Veränderlichen $x$ und $t$ betrachtet, in dem oben definirten
Sinne relativ prim sind, sich auch, als Functionen von $t$
allein betrachtet, und nach dem Algorithmus des grössten gemeinschaftlichen
Theilers behandelt (\S.\,6), als relativ prim erweisen
müssen. Denn wenn sie einen gemeinsamen Theiler hätten, der in
Bezug auf $t$ ganz, in Bezug auf die $x$ gebrochen wäre, so liesse
sich eine ganze Function $T$ von $x$ und $t$ und eine ganze Function
$P$ der $x$ allein so bestimmen, dass
\[
P\,F(t) = T\,F_1(t), \qquad P\,f(t) = T\,f_1(t)
\]
wäre, worin $F_1, f_1$ ganze Functionen ohne gemeinsamen Theiler
%CONT>>
